\chapter{1000 以内的数}\label{ch:g2:numbers}

一年级已学过把个数捆成十（见\cref{ch:g1:tens}）；现在，十个十
捆成一个\emph{百}，数的范围一直扩大到
$1000$。三个数字，各管一位：百位、十位、个位。

\section{百位、十位、个位}

\begin{definition}[三个数位]\label{def:g2:numbers:places}
一个\emph{百}就是十个十，即十捆，每捆十个。
数 $346$ 由三部分组成：
\[
346 = 3 \text{ 个百} + 4 \text{ 个十} + 6 \text{ 个一}
= 300 + 40 + 6 .
\]
\end{definition}

\begin{omfigure}
\begin{tikzpicture}[scale=0.85]
  \draw (0,0) grid[xstep=2.3, ystep=0.85] (6.9,1.7);
  \node[font=\small] at (1.15,1.275) {百位};
  \node[font=\small] at (3.45,1.275) {十位};
  \node[font=\small] at (5.75,1.275) {个位};
  \node[font=\large, omDef] at (1.15,0.425) {$3$};
  \node[font=\large, omDef] at (3.45,0.425) {$4$};
  \node[font=\large, omDef] at (5.75,0.425) {$6$};
\end{tikzpicture}

{\small 数 $346$ 的数位表，读作“三百四十六”。
（更大的数留待\cref{ch:g3:numbers}再学。）}
\end{omfigure}

\begin{example}[零占位]\label{ex:g2:numbers:zero}
“三百零五”就是 $305$：三个百，\emph{没有}
十位，五个一。$0$ 让 $5$ 留在个位上。没有它，
$35$ 就变成了另一个完全不同的数。
\end{example}

\begin{example}[数位之间的换算]\label{ex:g2:numbers:exchange}
十个一是一个十，十个十是一个百。所以 $10$ 个十
$= 100$，$2$ 个百 $= 20$ 个十。数 $250$ 含有
$25$ 个十，而不仅是数字 $5$。
\end{example}

\section{比较与排序}

\begin{method}[比较三位数]\label{met:g2:numbers:compare}
\begin{enumerate}
  \item 位数多的数大：$103 > 98$；
  \item 位数相同时 --- 从左往右比较：先比百位，再比
        十位，再比个位。第一个不同的数位决定大小：
        $472 > 468$，因为十位上 $7 > 6$。
\end{enumerate}
\end{method}

\begin{example}\label{ex:g2:numbers:order}
将 $530$、$305$、$503$ 排序：先比百位（$5$、$3$、$5$），
$305$ 排在最前；$530$ 与 $503$ 之间由十位决定（$3 > 0$）：
\[
305 < 503 < 530 .
\]
\end{example}

\section{数轴上的跳跃}

\begin{method}[加整十数与整百数]\label{met:g2:numbers:jumps}
在数轴上：
\begin{enumerate}
  \item $+1$ / $-1$：改变个位；
  \item $+10$ / $-10$：改变十位（$362 + 10 = 372$）；
  \item $+100$ / $-100$：改变百位（$362 + 100 = 462$）。
\end{enumerate}
每次只变动一个数位 --- 除非某一位“满十”，
并向前一位进位，例如 $195 + 10 = 205$。
\end{method}

\begin{omfigure}
\begin{tikzpicture}[scale=1.0]
  \draw[-Stealth] (-0.3,0) -- (10.6,0);
  \foreach \x/\l in {0.5/{362}, 5.5/{372}, 9.5/{472}}
    { \draw (\x,-0.1) -- (\x,0.1);
      \node[below, font=\small] at (\x,-0.12) {$\l$}; }
  \draw[-Stealth, omDef, thick] (0.5,0.28) .. controls (2.5,1.0) ..
    (5.5,0.28);
  \node[omDef, font=\small] at (3.0,1.05) {$+10$};
  \draw[-Stealth, omThm, thick] (5.5,0.28) .. controls (7.5,1.05) ..
    (9.5,0.28);
  \node[omThm, font=\small] at (7.5,1.1) {$+100$};
\end{tikzpicture}

{\small 从 $362$ 出发：小跳 $+10$ 改变十位，
大跳 $+100$ 改变百位。}
\end{omfigure}

\section{练习}

\begin{exercise}[$\star$]\label{exo:g2:numbers:1}
数一数下列各数中各有几个百、几个十和几个一：$258$；\ $470$；\ $600$；\
$909$？
\end{exercise}

\begin{exercise}[$\star$]\label{exo:g2:numbers:2}
写出下列各数：$4$ 个百、$0$ 个十、$7$ 个一；\ $6$ 个百
和 $3$ 个十；\ 十个百。
\end{exercise}

\begin{exercise}[$\star$]\label{exo:g2:numbers:3}
用数字写出：二百一十六；\ 五百零四；\
七百九十。
\end{exercise}

\begin{exercise}[$\star$]\label{exo:g2:numbers:4}
按\cref{def:g2:numbers:places}的方法分解：$438$；\ $702$；\
$560$。
\end{exercise}

\begin{exercise}[$\star$]\label{exo:g2:numbers:5}
抄写下列各式，并在空处填入 $<$ 或 $>$：
\[
327 \;?\; 372, \qquad
580 \;?\; 508, \qquad
199 \;?\; 200, \qquad
640 \;?\; 604 .
\]
\end{exercise}

\begin{exercise}[$\star$]\label{exo:g2:numbers:6}
从小到大排列：$430$；\ $304$；\ $340$；\ $403$。
\end{exercise}

\begin{exercise}[$\star$]\label{exo:g2:numbers:7}
用跳数线的方法计算：$253 + 10$；\quad $253 + 100$；
\quad $480 - 10$；\quad $716 - 100$。
\end{exercise}

\begin{exercise}[$\star$]\label{exo:g2:numbers:8}
接着数下去：$195, 205, 215, ?, ?$。\quad
$300, 400, 500, ?, ?$。\quad 第一题要细心哦！
\end{exercise}

\begin{exercise}[$\star$]\label{exo:g2:numbers:9}
一家商店有 $6$ 盒铅笔，每盒 $100$ 支，还有
$5$ 支散放的铅笔。一共有多少支铅笔？
\end{exercise}

\begin{exercise}[$\star\star$]\label{exo:g2:numbers:10}
用数字 $2$、$5$ 和 $8$ 各用一次，写出最大的
三位数和最小的三位数。它们相差多少？
（可以用数线跳着计算。）
\end{exercise}
